% Exact statement for production Rethlas and source audit, 26 August 2026.
% This proposition isolates the original Gubler--Kunnemann module category;
% it does not assert the unresolved continuous analytic Bedford--Taylor case.

\begin{proposition}
\label{prop:p21p3-delta-module}
Let \(K\) be an algebraically closed complete non-Archimedean field with
nontrivial absolute value, let \(X\) be an \(n\)-dimensional algebraic
\(K\)-variety, and let \(W\subset X^{\mathrm{an}}\) be open.  Write
\(B_{\mathrm{GK}}^{a,b}(W)\) for the Gubler--K\"unnemann delta-forms and
\(E_{\mathrm{GK}}^{p,q}(W)\) for their delta-currents.  Let
\[
 \alpha\in B_{\mathrm{GK}}^{a,b}(W),
 \qquad
 T\in E_{\mathrm{GK}}^{p,q}(W),
\]
where all bidegrees occurring below are in the range \(0,\ldots,n\).  Then
\begin{align}
 \mathrm d'(\alpha\wedge T)
 &=\mathrm d'\alpha\wedge T
   +(-1)^{a+b}\alpha\wedge\mathrm d'T,
 \label{eq:p21p3-dprime-module}\\
 \mathrm d''(\alpha\wedge T)
 &=\mathrm d''\alpha\wedge T
   +(-1)^{a+b}\alpha\wedge\mathrm d''T,
 \label{eq:p21p3-ddprime-module}
\end{align}
and consequently
\begin{align}
 \mathrm d'\mathrm d''(\alpha\wedge T)
 ={}&\mathrm d'\mathrm d''\alpha\wedge T
   +(-1)^{a+b+1}\mathrm d''\alpha\wedge\mathrm d'T
 \notag\\
 &+(-1)^{a+b}\mathrm d'\alpha\wedge\mathrm d''T
   +\alpha\wedge\mathrm d'\mathrm d''T.
 \label{eq:p21p3-dd-module}
\end{align}
In particular, if \(\mathrm d'T=\mathrm d''T=0\), then
\begin{equation}
 \mathrm d'\mathrm d''(\alpha\wedge T)
 =\mathrm d'\mathrm d''\alpha\wedge T.
 \label{eq:p21p3-closed-module}
\end{equation}
Thus form-level vanishing
\(\mathrm d'\mathrm d''\alpha=0\) in
\(B_{\mathrm{GK}}^{a+1,b+1}(W)\) forces
\(\mathrm d'\mathrm d''(\alpha\wedge T)=0\) as a delta-current.

Assume again that \(\mathrm d'T=\mathrm d''T=0\).  The same vanishing
conclusion follows from the functorial current-level hypothesis
\begin{equation}
 \bigl[f^*(\mathrm d'\mathrm d''\alpha)|_{W'}\bigr]_{D}=0
 \label{eq:p21p3-all-pullbacks}
\end{equation}
for every morphism \(f:X'\to X\) of algebraic \(K\)-varieties and every
open \(W'\subset (X')^{\mathrm{an}}\) with
\(f^{\mathrm{an}}(W')\subset W\), where the
left-hand side is the ordinary current on \(W'\) associated to the pulled
back generalized delta-form.

This statement supplies no multiplication of an arbitrary continuous or
piecewise-smooth function by an arbitrary delta-current.  Indeed,
\(B_{\mathrm{GK}}^{0,0}=A^{0,0}\), so a degree-zero multiplier in the
displayed module is smooth.  In this smooth degree-zero case there is a
useful sharper conclusion: if \(a=b=0\),
\(\mathrm d'T=\mathrm d''T=0\), and the single ordinary current
\begin{equation}
 \bigl[\mathrm d'\mathrm d''\alpha\bigr]_{D}=0
 \quad\text{on }W,
 \label{eq:p21p3-smooth-single-current}
\end{equation}
then \(\mathrm d'\mathrm d''(\alpha\wedge T)=0\).  By contrast, for a
general delta-form of general bidegree, vanishing of only its associated
ordinary current on \(W\) need not provide the form-level premise via the
original Gubler--K\"unnemann detection theorem; that theorem requires all
algebraic pullbacks.  Moreover, although the generalized piecewise-smooth
space \(P\!S\!P\) is an algebra and its elements define delta-currents, the
construction does not furnish a \(P\!S\!P\)-module structure on arbitrary
delta-currents extending the displayed \(B_{\mathrm{GK}}\)-module.  Hence
the proposition does not settle the continuous analytic Bedford--Taylor
question of
\cite[\S8.6.1 and Definition~8.6.2]{CLD25}.
No positivity, symmetry, or order hypothesis on \(T\) was used.
\end{proposition}
